\documentclass[lettersize,journal]{IEEEtran}
\usepackage{amsmath,amsfonts}
\usepackage{algorithmic}
\usepackage{algorithm}
\usepackage{array}
\usepackage[caption=false,font=normalsize,labelfont=sf,textfont=sf]{subfig}
\usepackage{textcomp}
\usepackage{stfloats}
\usepackage{url}
\usepackage{verbatim}
\usepackage{graphicx}
\usepackage{cite}
\usepackage{multirow,array}
\usepackage{amssymb}
\usepackage{cite}
\makeatletter
\newcommand{\rmnum}[1]{\romannumeral #1} % 小写罗马数字
\newcommand{\Rmnum}[1]{\expandafter\@slowromancap\romannumeral #1@} % 大写罗马数字

\usepackage{dblfloatfix}
\makeatother
\usepackage{float}
\usepackage{stfloats}
\usepackage{placeins}
\hyphenation{op-tical net-works semi-conduc-tor IEEE-Xplore}
% updated with editorial comments 8/9/2021

\begin{document}

\title{High-Confidence Pseudo-Label Generation and Learning for  Unsupervised Speaker Verification}

\author{Heng Wu, Yishuang Li, Hongchen Pan, Lin Li~\IEEEmembership{Member,~IEEE}, and Qingyang Hong~\IEEEmembership{Member,~IEEE}
        % <-this % stops a space
% \thanks{This paper was produced by the IEEE Publication Technology Group. They are in Piscataway, NJ.}% <-this % stops a space
% \thanks{Manuscript received April 19, 2021; revised August 16, 2021.}
}

% The paper headers
\markboth{IEEE TRANSACTIONS ON AUDIO, SPEECH, AND LANGUAGE PROCESSING,,~Vol.~1, No.~1, January ~2025}%
{Shell \MakeLowercase{\textit{et al.}}: A Sample Article Using IEEEtran.cls for IEEE Journals}

%\IEEEpubid{0000--0000/00\$00.00~\copyright~2021 IEEE}
% Remember, if you use this you must call \IEEEpubidadjcol in the second
% column for its text to clear the IEEEpubid mark.

\maketitle

\begin{abstract}
Unsupervised speaker verification can exploit large-scale unlabeled
speech without speaker identity annotations, but iterative training
remains sensitive to unreliable pseudo-labels. In particular,
uncertain cluster assignments and category fragmentation can degrade
supervision quality and propagate errors across training iterations.
To address these issues, we propose a High-Confidence Pseudo-Label
(HCPL) generation and learning framework for unsupervised speaker
verification.
HCPL generates two clustering assignments using different speech
pre-trained models and exploits cross-model assignment agreement to
identify high-confidence pseudo-labels. A pseudo-label merging
strategy is subsequently applied to reduce fragmentation among highly
similar clusters. During speaker encoder training, a dynamic label
noise correction objective combines pseudo-label supervision with
confidence-weighted model predictions to mitigate residual label
noise. Meanwhile, samples excluded from pseudo-label-guided training
are retained for contrastive learning with hard negative mining.
Experiments using VoxCeleb2 for training show that the proposed
framework reaches its best reported performance after three
refinement iterations, achieving equal error rates of 1.01\%,
1.36\%, and 2.53\% on Vox1-O, Vox1-E, and Vox1-H, respectively.
The results demonstrate the effectiveness of improving pseudo-label
reliability at both the generation and optimization stages.
Importantly, the proposed unsupervised system achieves performance
close to that of the fully supervised system on Vox1-O, without using
speaker identity annotations during training.
\end{abstract}

\begin{IEEEkeywords}
Unsupervised speaker verification, pseudo-label, label noise correction, hard negative mining, speech pre-trained models.
\end{IEEEkeywords}

\section{Introduction}
\IEEEPARstart{S}{peaker} verification determines whether a test
utterance belongs to a claimed speaker and is widely used in
speech-based authentication. Deep neural network-based systems
[1]--[5] have substantially improved speaker verification
performance over traditional approaches [6], [7]. These advances, however, generally depend on
large-scale speaker-labeled corpora. Although speech data can be
collected at scale, reliable speaker identities are often unavailable
or require costly manual annotation and additional identity-related
supervision. This mismatch between the availability of speech data
and speaker labels motivates the development of unsupervised speaker
verification methods.

Existing unsupervised speaker verification studies mainly learn
speaker representations through self-supervised objectives or
iteratively train speaker encoders using clustering-derived
pseudo-labels. Contrastive learning methods inspired by SimCLR [8]
and MoCo [9] construct positive and negative speech pairs to learn
speaker-discriminative representations without manual labels
[10]--[14]. Iterative frameworks further alternate between
representation learning and pseudo-label generation, allowing the
speaker encoder and clustering assignments to be progressively
updated [15]--[20]. These methods have narrowed the gap between
unsupervised and supervised systems, but their performance remains
sensitive to the quality of the supervision generated from unlabeled
data.

A major difficulty lies in the reliability of pseudo-labels,
particularly during the initial iterations. Contrastive objectives
may treat utterances from the same speaker as negative pairs, thereby
introducing false negatives [21]. In clustering-based frameworks,
samples near ambiguous cluster boundaries may receive incorrect
assignments, while utterances from one speaker may be divided into
multiple pseudo-label classes. Once used as supervision, these errors
can be reinforced in subsequent iterations. Consequently, low-
confidence assignments, category fragmentation, and cross-iteration
error propagation jointly limit the stability and effectiveness of
iterative unsupervised training.

Previous studies have addressed parts of this problem from different
perspectives. Self-distillation methods [22]--[30] improve
representation stability without relying on explicit negative pairs,
and speech pre-trained models [31]--[35] provide transferable
representations for clustering and downstream speaker modeling.
PTM-based embeddings have also been shown to improve first-round
pseudo-label generation [36]. In parallel, noise-resistant learning
methods use regularized objectives, confidence-aware optimization,
classifier adjustment, and consistency-based sample selection
[37]--[40]. Other approaches refine clustering outputs through
confidence filtering, category pruning, GMM-based sample selection,
or iterative label correction [36], [41], [42]. These methods improve
either representation quality, training robustness, or pseudo-label
selection, but the reliability of initial assignments, category
fragmentation, and residual label noise are often handled separately.

This work approaches the problem from both pseudo-label generation
and pseudo-label-guided learning. Our objective is to construct a
more reliable initial supervisory set, reduce obvious category
fragmentation, and limit the influence of remaining label noise
during iterative optimization. To this end, we propose a
High-Confidence Pseudo-Label framework, termed HCPL. Two speech
pre-trained models are first used to generate independently obtained
clustering assignments. Since the models employ different
pre-training objectives and may produce non-identical clustering
errors, cross-model assignment consistency is used as a reliability
cue to retain high-confidence samples. A pseudo-label merging
procedure is then applied to consolidate highly similar clusters.

The filtered and merged pseudo-labels are used to train the speaker
encoder with a reliability-aware objective. A Dynamic Label Noise
Correction loss combines pseudo-label supervision with
confidence-weighted model predictions, providing soft correction for
residual label noise. The samples excluded by high-confidence
filtering are not discarded; instead, they are utilized through
contrastive learning with hard negative mining. The resulting
framework therefore separates the use of unlabeled data according to
supervision reliability while progressively refining the
pseudo-labels across iterations.

The contributions of this work are summarized as follows:

\begin{enumerate}
    \item We introduce an HCPL generation procedure that combines
    dual-PTM assignment agreement with pseudo-label merging. The
    former identifies a high-confidence supervisory subset, while
    the latter reduces category fragmentation among highly similar
    clusters.

    \item We develop a training strategy that combines dynamic
    label-noise correction with hard-negative-based contrastive
    learning. The strategy uses confidence-aware soft correction for
    pseudo-labeled samples and retains the filtered samples for
    self-supervised representation learning.

    \item Experiments on VoxCeleb2 show that the proposed framework
    improves first-round pseudo-label quality and speaker
    verification performance. The complete system reaches its best
    reported performance within three iterations and achieves EERs
    of 1.01\%, 1.36\%, and 2.53\% on Vox1-O, Vox1-E, and Vox1-H,
    respectively.
\end{enumerate}


The remainder of this paper is organized as follows. Section \uppercase\expandafter{\romannumeral2} reviews related work on unsupervised speaker verification and noise-resistant learning. Section \uppercase\expandafter{\romannumeral3} presents the proposed HCPL framework for high-confidence pseudo-label generation and refinement. Section \uppercase\expandafter{\romannumeral4} introduces the proposed joint optimization strategy for iterative training. Section \uppercase\expandafter{\romannumeral5} describes the experimental setup and reports the corresponding results and analyses. Finally, Section \uppercase\expandafter{\romannumeral6} concludes this paper.

\section{Related Work}

\subsection{Contrastive Learning for Speaker Representation}

Contrastive learning has been widely adopted for unsupervised
speaker representation learning because it can exploit unlabeled
speech through pairwise similarity constraints. Inspired by
frameworks such as SimCLR~[8] and MoCo~[9], previous studies
[10]--[14] construct positive pairs from augmented views or segments
of the same utterance and generally treat samples from different
utterances as negative pairs. The resulting objectives encourage
representations from the same utterance to remain close while
separating those from different utterances.

The effectiveness of contrastive learning, however, depends on the
validity of the negative-pair assumption. In large-scale unlabeled
speech corpora, different utterances may originate from the same
speaker and can therefore be incorrectly treated as negatives. Such
false negatives introduce repulsive constraints between
same-speaker samples and may degrade the learned representation
[21]. Although existing studies improve contrastive speaker learning
through data augmentation, margin-based objectives, positive-pair
construction, and sample selection [10]--[14], contrastive learning
alone does not directly ensure the reliability of clustering-derived
pseudo-labels in iterative systems.

\subsection{Self-Distillation for Speaker Representation Learning}

Self-distillation provides an alternative self-supervised paradigm
that avoids explicit negative-pair construction. Methods based on
teacher--student architectures enforce prediction consistency across
multiple augmented views and have been applied to unsupervised
speaker representation learning [22]--[26]. By removing manually
defined negative pairs, these methods reduce the direct influence of
false-negative sampling and improve representation stability.

Self-distillation has also been incorporated into iterative
unsupervised speaker verification frameworks [27]--[30]. In these
systems, the learned representations are clustered to generate
pseudo-labels for subsequent model refinement. Although
self-distillation improves the robustness of the underlying
representations, it does not directly guarantee reliable clustering
assignments. Ambiguous samples may still receive incorrect
pseudo-labels, and utterances from the same speaker may be separated
into multiple clusters. Additional pseudo-label filtering and
refinement mechanisms are therefore required during iterative
training.

\subsection{Pre-Trained Models for Unsupervised Speaker Verification}

Speech pre-trained models have become important feature extractors
for unsupervised speaker verification. Models such as wav2vec
2.0~[32], HuBERT~[33], UniSpeech-SAT~[34], and WavLM~[35] learn
transferable representations from large-scale unlabeled speech and
encode acoustic, phonetic, and speaker-related information that can
support downstream speech tasks [44].

Previous studies [36], [49] have shown that PTM-based features can
improve clustering initialization and unsupervised speaker
verification performance. In particular, PTM representations can
provide informative embeddings during the early stage of iterative
pseudo-label learning and can be combined with downstream speaker
encoders to incorporate both pre-trained and task-specific
information.

Nevertheless, improved representations do not necessarily produce
fully reliable pseudo-labels. Early clustering results may still
contain uncertain assignments and fragmented speaker categories.
Moreover, PTMs trained with different objectives may produce
non-identical clustering partitions. These properties motivate the
use of cross-model assignment consistency as a reliability cue,
rather than relying on the clustering output of a single PTM.

\subsection{Iterative Frameworks for Unsupervised Speaker Verification}

Iterative optimization has been widely used to reduce the performance
gap between unsupervised and supervised speaker verification
systems [15]--[20]. A typical framework alternates between two
stages. First, speaker embeddings are extracted from unlabeled speech
and clustered to generate pseudo-labels. Second, the pseudo-labels
are used to train or refine a speaker encoder. The updated encoder
then produces new embeddings for the next round of clustering and
model optimization.

This iterative procedure enables speaker representations and
pseudo-labels to improve jointly. Its effectiveness, however, is
sensitive to the quality of the initial clustering results.
Incorrect assignments generated in early iterations may be
reinforced in later rounds, while category fragmentation can provide
inconsistent supervision for utterances from the same speaker.
Consequently, improvements in representation learning alone may not
be sufficient when the generated pseudo-labels remain unreliable.

The proposed method follows the iterative unsupervised speaker
verification setting but focuses on controlling pseudo-label
reliability at different stages. Cross-model agreement is first used
to identify a high-confidence pseudo-label subset, and pseudo-label
merging is applied to reduce category fragmentation. During model
training, dynamic label-noise correction mitigates residual
supervision noise, while contrastive learning with hard negative
mining utilizes the samples excluded from pseudo-label-guided
optimization. In this way, the framework combines initial
pseudo-label selection, category refinement, and noise-aware learning
within the iterative pipeline.


\begin{figure*}[!t]
\centering
\includegraphics[width=7in]{v3.png}
\caption{\centering{The overall architecture of the proposed method.
Step 1–4: High-confidence pseudo-label (HCPL) generation via dual PTMs.
Step 5–6: Model training with hard negative mining-based contrastive loss and dynamic label correction loss. }}
\label{fig:frame}
\end{figure*}

\section{High-Confidence Pseudo-Label Generation Framework}

Existing iterative unsupervised speaker verification systems are fundamentally constrained by unreliable pseudo-label supervision during early training iterations. In practice, pseudo-labels generated from clustering often suffer from two critical issues: low-confidence assignments caused by noisy clustering boundaries and category fragmentation in which samples from the same speaker are partitioned into multiple clusters. These errors are repeatedly propagated throughout iterative optimization, resulting in unstable convergence and degraded verification performance.

To address these limitations, we propose the HCPL framework based on dual PTMs. The proposed framework aims to improve pseudo-label reliability before downstream model optimization by jointly suppressing noisy label assignments and reducing category fragmentation. As illustrated in Fig.~\ref{fig:frame}, HCPL consists of four stages:
1) extracting speaker embeddings using PTMs;
2) generating initial pseudo-labels through clustering;
3) performing cross-model consistency filtering to retain high-confidence pseudo-labels; and
4) merging fragmented categories to obtain a more compact and reliable pseudo-label set.

\subsection{Pre-Trained Models for Reliable Pseudo-Labeling}

PTMs have recently demonstrated strong capability in learning transferable speech representations from large-scale unlabeled corpora. Prior studies [50] show that speaker-related information is unevenly distributed across Transformer layers: lower layers mainly encode local acoustic and channel characteristics, whereas higher layers are increasingly dominated by linguistic and semantic information. Consequently, intermediate Transformer layers often provide more discriminative speaker representations for speaker verification tasks. Following this observation, we extract speaker embeddings from the $L$-th Transformer layer, where $L$ is selected empirically according to validation performance.

In this work, we employ several representative speech PTMs, including wav2vec~2.0, HuBERT, UniSpeech-SAT, and WavLM. Although these models share similar encoder architectures, their distinct pre-training objectives lead to complementary representation characteristics and different clustering error distributions. We argue that such cross-model diversity provides a natural reliability cue for pseudo-label generation: pseudo-label assignments that remain consistent across different PTMs are more likely to correspond to stable speaker structures rather than model-specific clustering noise. Based on this observation, we construct a dual-PTM framework in which cross-model agreement is utilized as a criterion for high-confidence pseudo-label selection.

\subsection{Pseudo-Label Generation via Infomap Clustering}

After extracting speaker embeddings from PTMs, clustering is performed to generate pseudo-labels for unlabeled speech samples. Traditional clustering methods such as K-means [51] mainly operate in Euclidean space and often struggle to capture the complex manifold structure of high-dimensional speech embeddings. To better model the intrinsic neighborhood relationships among speaker representations, we adopt the graph-based Infomap clustering algorithm [52], which has demonstrated strong effectiveness in speaker clustering tasks [36], [53].

Infomap identifies community structures by minimizing the expected coding length of information flow over a graph. In our implementation, each speaker embedding is treated as a graph node, while cosine similarities between embeddings define weighted graph edges. Random walks are then performed on the graph to model information flow between nodes. The clustering objective is formulated through the map equation:
\begin{equation}
L(M)
=
P_{\mathrm{exit}}L_{\mathrm{exit}}
+
\sum_{i=1}^{K}
P_iL_i,
\end{equation}
where $P_{\mathrm{exit}}$ denotes the probability of inter-cluster transitions, $L_{\mathrm{exit}}$ represents the average coding length for inter-cluster movement, and $P_i$ and $L_i$ denote the visitation probability and average coding length within the $i$-th cluster, respectively. During optimization, nodes are iteratively assigned to neighboring clusters that minimize the overall coding length until convergence is reached.

Compared with conventional clustering methods, Infomap can better capture the non-linear structural relationships among speaker embeddings, thereby producing more compact and semantically consistent pseudo-label clusters.

\subsection{High-Confidence Pseudo-Label Filtering}

Pseudo-label quality in iterative unsupervised speaker verification is highly sensitive to clustering errors during early iterations. In particular, unreliable pseudo-label assignments may introduce incorrect supervision signals that are repeatedly amplified throughout iterative optimization. To alleviate this issue, we propose a high-confidence filtering (HCF) strategy based on cross-model agreement between two PTMs.

The proposed method is motivated by the observation that different PTMs capture complementary speaker characteristics and exhibit distinct clustering error distributions. Consequently, pseudo-label assignments that remain consistent across multiple PTMs are more likely to reflect intrinsic speaker structures rather than model-specific noise. We therefore treat cross-model consistency as a reliability indicator for pseudo-label selection.

Specifically, let $Y_1$ and $Y_2$ denote two pseudo-label sets obtained from clustering the embeddings extracted by two different PTMs. The corresponding cluster sets are defined as
\begin{equation}
C_1=\{c_1^1,c_2^1,\ldots,c_m^1\},
\end{equation}
\begin{equation}
C_2=\{c_1^2,c_2^2,\ldots,c_n^2\},
\end{equation}
where $c_i^1$ and $c_j^2$ denote the cluster identities in $Y_1$ and $Y_2$, respectively.

Since clustering labels are permutation-invariant, the specific numerical values of cluster indices are arbitrary as long as the underlying sample grouping remains unchanged [54]. Based on this property, we first compute the overlap score between clusters:
\begin{equation}
S(i,j)
=
\left|
\left\{
x_k
\mid
x_k\in U_0,
Y_1(x_k)=c_i^1,
Y_2(x_k)=c_j^2
\right\}
\right|,
\end{equation}
where $U_0$ denotes the unlabeled dataset and $x_k$ represents the $k$-th speech sample.

The Hungarian algorithm [55] is then employed to obtain the optimal cluster alignment:
\begin{equation}
\mathrm{Match}
=
\arg\max_{\pi}
\sum_{(i,j)\in\pi}
S(i,j),
\end{equation}
where $\pi$ denotes a feasible cluster matching configuration. Based on the obtained matching relationship, the pseudo-labels in $Y_2$ are aligned with those in $Y_1$:
\begin{equation}
\mathrm{Mapping}(c_j^2)=c_i^1,
\quad
\forall(i,j)\in\mathrm{Match},
\end{equation}
and the aligned pseudo-labels are defined as
\begin{equation}
Y_2^{\mathrm{aligned}}(x_k)
=
\mathrm{Mapping}(Y_2(x_k)).
\end{equation}

Based on the aligned pseudo-label assignments, the high-confidence sample set is defined as
\begin{equation}
\mathcal{I}
=
\left\{
x_k
\mid
Y_1(x_k)
=
Y_2^{\mathrm{aligned}}(x_k)
\right\}.
\end{equation}

The final high-confidence pseudo-label set is therefore obtained as
\begin{equation}
Y_f
=
Y_1(\mathcal{I}).
\end{equation}

By filtering pseudo-labels before downstream optimization, the proposed HCF strategy reduces error propagation during iterative training and provides a more stable initialization for subsequent model refinement.

\subsection{Pseudo-Label Merging}

Although high-confidence filtering removes a substantial portion of
unreliable pseudo-label assignments, category fragmentation may still
remain. Specifically, utterances from the same speaker may be divided
into multiple clusters because of acoustic variability, recording
conditions, and imperfect clustering boundaries. Such over-segmentation
increases the number of pseudo-label classes and may introduce
inconsistent supervision during subsequent model training.

To alleviate this problem, we apply a centroid-based pseudo-label
merging strategy to the high-confidence pseudo-label set. For the
$i$-th pseudo-label cluster, its centroid is computed as
\begin{equation}
\mathbf{c}_i
=
\frac{1}{N_i}
\sum_{k=1}^{N_i}
\mathbf{e}_{i,k},
\end{equation}
where $N_i$ denotes the number of samples in the $i$-th cluster and
$\mathbf{e}_{i,k}$ represents the embedding of its $k$-th sample.

The similarity between clusters $i$ and $j$ is measured by the cosine
similarity between their centroids:
\begin{equation}
s_{i,j}
=
\cos(\mathbf{c}_i,\mathbf{c}_j).
\end{equation}

Rather than using a fixed absolute threshold, we define the merging
threshold relative to the maximum inter-cluster similarity in the
current pseudo-label set:
\begin{equation}
\theta_{\mathrm{PLM}}
=
\mu
\max_{i\neq j}
s_{i,j},
\end{equation}
where $\mu\in[0,1]$ is a scaling coefficient controlling the
strictness of the merging criterion. A larger value of $\mu$ results
in more conservative merging, whereas a smaller value allows a larger
number of cluster pairs to satisfy the merging condition. In all
experiments, $\mu$ is empirically fixed to $0.97$ and kept unchanged
throughout the training procedure.

Two clusters are merged when their centroid similarity satisfies
\begin{equation}
s_{i,j}
\geq
\theta_{\mathrm{PLM}}.
\end{equation}

The relative threshold adapts the merging criterion to the similarity
scale of the current embedding space, avoiding the need for a
dataset-independent absolute similarity threshold. The relatively high
value of $\mu$ imposes a conservative merging condition, under which
only cluster pairs whose centroid similarities are close to the
maximum observed similarity are considered for consolidation. This
design aims to reduce obvious category fragmentation while limiting
the risk of merging clusters from different speakers.

After a pair of clusters is merged, the corresponding samples are
assigned the same pseudo-label. The resulting merged label set,
denoted as $\tilde{Y}$, is then used as the initial supervision for
downstream speaker verification training.
\begin{algorithm}[t]
\caption{\textbf{HCPL:} High-Confidence Pseudo-Label Generation Framework}
\label{alg:alg1}
\begin{algorithmic}
\STATE
\STATE \textbf{Input:} Unlabeled audio set $U_0$, PTM-A, PTM-B, clustering algorithm $\mathcal{C}(\cdot)$, merging threshold $\theta_{\mathrm{PLM}}$
\STATE \textbf{Output:} Reliable pseudo-label set $\tilde{Y}$

\STATE \textbf{Step 1: Speaker embedding extraction}
\FOR{each $u \in U_0$}
    \STATE $\mathbf{E}_A(u)\gets \mathrm{PTM\mbox{-}A}(u)$
    \STATE $\mathbf{E}_B(u)\gets \mathrm{PTM\mbox{-}B}(u)$
\ENDFOR

\STATE \textbf{Step 2: Pseudo-label generation}
\STATE $\{Y_A(u)\}_{u\in U_0}\gets \mathcal{C}(\{\mathbf{E}_A(u)\})$
\STATE $\{Y_B(u)\}_{u\in U_0}\gets \mathcal{C}(\{\mathbf{E}_B(u)\})$

\STATE \textbf{Step 3: High-confidence filtering (HCF)}
\STATE Construct matching matrix between $Y_A$ and $Y_B$
\STATE $(\pi_A,\pi_B)\gets \mathrm{Hungarian}(\cdot)$
\STATE $U^{\mathrm{hc}}\gets\varnothing$

\FOR{each $u\in U_0$}
    \IF{$\pi_A(Y_A(u))=\pi_B(Y_B(u))$}
        \STATE $U_{\mathrm{hc}}\gets U_{\mathrm{hc}}\cup\{u\}$
    \ENDIF
\ENDFOR

\STATE \textbf{Step 4: Pseudo-label merging (PLM)}
\STATE Compute cluster centroids $\{\mathbf{c}_k\}$ for samples in $U_{\mathrm{hc}}$

\FOR{each pair $(\mathbf{c}_i,\mathbf{c}_j)$}
    \STATE $\mathrm{sim}(i,j)\gets \mathrm{CosineSim}(\mathbf{c}_i,\mathbf{c}_j)$
    \IF{$\mathrm{sim}(i,j)\ge\theta_{\mathrm{PLM}}$}
        \STATE Merge clusters $i$ and $j$
    \ENDIF
\ENDFOR

\STATE $\tilde{Y}\gets$ merged pseudo-labels of samples in $U^{\mathrm{hc}}$

\STATE \textbf{return} $\tilde{Y}$

\end{algorithmic}
\end{algorithm}
\section{Learning with Noisy Pseudo-Labels}

After generating reliable pseudo-labels through the proposed HCPL framework, the remaining challenge lies in mitigating residual pseudo-label noise during downstream optimization. Although high-confidence filtering substantially improves pseudo-label reliability, incorrect assignments may still exist and can negatively affect iterative model refinement. To address this issue, we design a unified training framework that jointly integrates PTM-enhanced feature fusion, Dynamic Label Noise Correction (DLNC), and contrastive learning with hard negative mining (HNM). The proposed framework aims to improve supervision reliability while fully exploiting unlabeled speech data for discriminative speaker representation learning.

\subsection{Speaker Verification Backbone}

Recent studies [49] have shown that integrating PTM-based representations with conventional acoustic features can improve speaker verification performance by leveraging complementary information across different feature spaces. Motivated by this observation, we adopt a hybrid feature extraction framework that combines Transformer-based PTM representations with conventional FBank features.

Specifically, multi-layer Transformer outputs from the PTM are adaptively fused to capture speaker-related information distributed across different semantic levels. Meanwhile, FBank features are processed through a convolutional neural network (CNN) encoder to preserve local acoustic characteristics. The fused representations are subsequently fed into an ECAPA-TDNN backbone [7] for downstream speaker embedding extraction and model optimization.

\subsection{Dynamic Label Noise Correction}

Although the proposed HCPL framework improves pseudo-label
quality through high-confidence filtering and pseudo-label
merging, residual errors may still remain because of imperfect
clustering boundaries and speaker ambiguity. Directly treating
all pseudo-labels as equally reliable may introduce biased
supervision and reinforce incorrect assignments during iterative
training. To mitigate this problem, we propose a Dynamic Label
Noise Correction (DLNC) strategy that performs confidence-aware
soft correction at the objective level.

Unlike hard label correction methods that directly replace a
pseudo-label with the current model prediction, DLNC preserves
the original pseudo-label as a stable supervision anchor and
introduces the model prediction as an auxiliary correction signal.
The contribution of this correction signal is dynamically
controlled by both the global training state and the uncertainty
of each individual sample.

For a sample $x_i$, let
\begin{equation}
P_{i,k}=p(k\mid x_i;\Theta)
\end{equation}
denote the posterior probability assigned by the current model
to class $k$, where $\Theta$ represents the model parameters.
The model-predicted label is defined as
\begin{equation}
\hat{y}_i=\arg\max_k P_{i,k}.
\end{equation}

\subsubsection{Global Prediction Reliability}

The reliability of model predictions changes throughout
optimization. At the beginning of training, model predictions
are generally less stable and should therefore have only a
limited influence on pseudo-label correction. As training
progresses, the learned representation becomes more
discriminative, and high-confidence model predictions can
provide increasingly informative correction signals.

To characterize the overall prediction confidence at training
step $t$, we compute the average maximum posterior probability
within a mini-batch:
\begin{equation}
\hat{P}_t
=
\frac{1}{B}
\sum_{i=1}^{B}
P_{i,\hat{y}_i},
\end{equation}
where $B$ denotes the mini-batch size. Since the batch-level
estimate may fluctuate across training steps, we apply an
exponential moving average:
\begin{equation}
\delta_t
=
v\delta_{t-1}
+
(1-v)\hat{P}_t,
\end{equation}
where $v\in[0,1)$ is the momentum coefficient. The variable
$\delta_t$ provides a smoothed estimate of the global prediction
confidence and controls the overall strength of model-driven
self-correction.

\subsubsection{Sample-Level Prediction Reliability}

In addition to the global training state, individual samples
exhibit different levels of prediction uncertainty. A sample
whose largest posterior probability is only slightly higher than
the second-largest probability should contribute less to label
correction than a sample with a clearly separated prediction.

Let
\begin{equation}
P_{i,\mathrm{2nd}}
=
\max_{k\neq\hat{y}_i}P_{i,k}
\end{equation}
denote the second-largest posterior probability for sample
$x_i$. The posterior margin is defined as
\begin{equation}
m_i
=
P_{i,\hat{y}_i}
-
P_{i,\mathrm{2nd}}.
\end{equation}

We further normalize the margin by the maximum posterior
probability:
\begin{equation}
\omega_i
=
\frac{m_i}{P_{i,\hat{y}_i}}
=
1-
\frac{P_{i,\mathrm{2nd}}}{P_{i,\hat{y}_i}}.
\end{equation}

Since $P_{i,\hat{y}_i}\geq P_{i,\mathrm{2nd}}\geq 0$, the
reliability score satisfies $\omega_i\in[0,1]$. A larger
$\omega_i$ indicates that the predicted class is more clearly
separated from competing classes and can therefore provide a
stronger correction signal.

\subsubsection{Confidence-Aware Soft Correction Objective}

Given the pseudo-label $y_i$, conventional pseudo-label
supervision is formulated as
\begin{equation}
\mathcal{L}_{\mathrm{PL}}
=
-\frac{1}{B}
\sum_{i=1}^{B}
\log P_{i,y_i}.
\end{equation}

This objective assumes that all pseudo-labels are correct and
equally reliable, making it sensitive to residual clustering
errors. To dynamically regulate the contribution of the
model-predicted supervision, we combine the global confidence
indicator $\delta_t$ and the sample-level reliability score
$\omega_i$:
\begin{equation}
\lambda_i
=
\delta_t\omega_i,
\end{equation}
where $\lambda_i\in[0,1]$ is the adaptive correction coefficient
for sample $x_i$.

The proposed DLNC objective is defined as
\begin{equation}
\mathcal{L}_{\mathrm{DLNC}}
=
-\frac{1}{B}
\sum_{i=1}^{B}
\left[
\log P_{i,y_i}
+
\lambda_i\log P_{i,\hat{y}_i}
\right].
\end{equation}

The first term retains the original pseudo-label as the primary
supervision signal, thereby preventing unstable model
predictions from immediately overwriting the clustering
assignment. The second term introduces a model-driven
self-correction signal whose contribution depends on both the
training stage and the uncertainty of the individual sample.

When the predicted label agrees with the pseudo-label,
i.e., $\hat{y}_i=y_i$, the sample-wise objective becomes
\begin{equation}
\mathcal{L}_i
=
-(1+\lambda_i)\log P_{i,y_i}.
\end{equation}
In this case, the two supervision sources reinforce the same
target, assigning a larger optimization weight to samples that
are supported by both clustering and model prediction.

When $\hat{y}_i\neq y_i$, the pseudo-label remains the primary
supervision anchor, while the model prediction provides an
auxiliary correction direction weighted by $\lambda_i$. During
the early stage of training, $\delta_t$ is relatively small, so
potentially unreliable predictions have only a limited effect.
As training progresses and the overall prediction confidence
increases, high-margin predictions receive larger correction
weights, allowing the model to gradually reduce the influence
of potentially incorrect pseudo-label supervision.

Therefore, DLNC simultaneously performs confidence-aware
soft target correction and reliability-based sample reweighting.
This gradual correction mechanism avoids abrupt label
replacement and reduces the risk of confirmation bias caused by
unreliable early-stage predictions.

\subsubsection{Reliability-Enhanced Margin Learning}

For samples whose model predictions agree with their pseudo-labels,
i.e., $\hat{y}_i=y_i$, we follow the reliability-aware margin
adjustment strategy introduced in [39]. Such samples are treated as
relatively reliable and are optimized with an adjusted AAM-Softmax
objective.

Specifically, the posterior probability of the target class is
computed as
\begin{equation}
P'_{i,\hat y_i}
=
\frac{
e^{s\cos(\theta_i^{\hat y_i}+\alpha)}
}
{
e^{s\cos(\theta_i^{\hat y_i}+\alpha)}
+
\sum_{j\neq \hat y_i}
e^{s(\cos(\theta_i^j)+r_{i,j})}
},
\end{equation}
where $s$ and $\alpha$ denote the scaling factor and angular margin,
respectively. Following [39], the adaptive margin for a non-target
class $j$ is defined as
\begin{equation}
r_{i,j}
=
\frac{1}{10}
\left(
1-\cos(\theta_i^j)
\right),
\qquad r_{i,j}\geq 0.
\end{equation}

This margin adjustment increases the separation constraint between
the target class and competing non-target classes for samples with
consistent pseudo-label and model prediction. In our framework, this
strategy is used as an auxiliary optimization component rather than
as a newly proposed contribution.

\begin{algorithm}[t]
\caption{Iterative Unsupervised Speaker Verification Training with HCPL Initialization}
\label{alg:alg2}
\begin{algorithmic}
\STATE 
\STATE \textbf{Input:} Unlabeled audio set $U_{0}$, pretrained encoders PTM-A and PTM-B, maximum iterations $K$, stopping criterion $\varepsilon$
\STATE \textbf{Output:} Trained speaker encoder $f_{\theta}^{(*)}$

\STATE 
\STATE \textbf{Initialization:} Randomly initialize $f_{\theta}^{(0)}$ or adopt a pretrained checkpoint

%--------------------------------------------------------
\STATE 
\STATE \textbf{Iteration 1: HCPL-based pseudo-labeling}
\STATE $(\tilde{Y}^{(1)}, U_{\mathrm{hc}}^{(1)})+U_{\mathrm{nl}}^{(1)} \gets \mathrm{HCPL}(U_{0})$ \hfill \COMMENT{Call Alg.~\ref{alg:alg1}, return high-confidence labels and discarded set}
\STATE Train $f_{\theta}^{(1)}$ on $(U_{\mathrm{hc}}^{(1)}, \tilde{Y}^{(1)}) + U_{\mathrm{nl}}^{(1)}$ with loss
\STATE \hspace{0.8cm} $\mathcal{L} = \mathcal{L}_{\mathrm{DLNC}} + \mathcal{L}_{\mathrm{HNM}}$
\STATE Evaluate performance metric $\mathcal{M}^{(1)}$

%--------------------------------------------------------
\STATE 
\STATE \textbf{Subsequent iterations: prediction-based pseudo labels}
\FOR{$k = 2$ to $K$}
    \STATE \textbf{Step 1: Generate pseudo labels from previous model}
    \FOR{each $u \in U_{0}$}
        \STATE $\hat{y}^{(k)}(u) \gets \mathrm{Predict}\big(f_{\theta}^{(k-1)}, u\big)$
    \ENDFOR
    \STATE $\hat{Y}^{(k)} \gets \{\hat{y}^{(k)}(u)\}_{u \in U_{0}}$
    
    \STATE 
    \STATE \textbf{Step 2: Supervised training with $\mathcal{L}_{\mathrm{DLNC}}$}
    \STATE Train $f_{\theta}^{(k)}$ on $(U_{0}, \hat{Y}^{(k)})$ with loss
    \STATE \hspace{0.8cm} $\mathcal{L} = \mathcal{L}_{\mathrm{DLNC}}$
    
    \STATE 
    \STATE \textbf{Step 3: Convergence check}
    \STATE Evaluate performance metric $\mathcal{M}^{(k)}$
    \IF{$|\mathcal{M}^{(k)} - \mathcal{M}^{(k-1)}| < \varepsilon$}
        \STATE \textbf{break}
    \ENDIF
\ENDFOR

\STATE 
\STATE $f_{\theta}^{(*)} \gets f_{\theta}^{(k)}$
\STATE \textbf{return} $f_{\theta}^{(*)}$

\end{algorithmic}
\end{algorithm}

\subsection{Contrastive Learning with Hard Negative Mining}

Although the proposed HCPL framework filters unreliable pseudo-labels, a portion of unlabeled data remains excluded from supervised optimization. Directly discarding these samples may result in inefficient utilization of unlabeled speech data. To address this issue, we introduce a self-supervised contrastive learning objective with hard negative mining (HNM) to further exploit the remaining unlabeled samples.

The proposed strategy aims to improve inter-speaker discrimination by focusing optimization on acoustically similar yet semantically different speaker representations, which are particularly challenging in large-scale speaker verification scenarios.

Let $U_{\mathrm{nl}}$ denote the unlabeled data excluded by high-confidence filtering. For the $i$-th audio sample in $U_{\mathrm{nl}}$, two non-overlapping 2-second speech segments are randomly sampled and fed into the speaker verification model to obtain embeddings $\mathbf{z}_i^1$ and $\mathbf{z}_i^2$. The two embeddings extracted from the same utterance form a positive pair, while embeddings originating from different utterances constitute negative pairs.

After $L_2$ normalization,
\begin{equation}
\bar{\mathbf{z}}
=
\frac{\mathbf{z}}{\|\mathbf{z}\|},
\end{equation}
the cosine similarity between two embeddings is computed as
\begin{equation}
S_{ij}
=
\bar{\mathbf{z}}_i^1
\cdot
\bar{\mathbf{z}}_j^2.
\end{equation}

The similarity matrix within a mini-batch is therefore represented as
\begin{equation}
S=
\begin{bmatrix}
S_{11} & S_{12} & \cdots & S_{1N} \\
S_{21} & S_{22} & \cdots & S_{2N} \\
\vdots & \vdots & \ddots & \vdots \\
S_{N1} & S_{N2} & \cdots & S_{NN}
\end{bmatrix},
\end{equation}
where $N$ denotes the mini-batch size.

Following a linear transformation
\begin{equation}
S'_{ij}
=
\omega S_{ij}+b,
\end{equation}
the contrastive loss is formulated as
\begin{equation}
\mathcal{L}_{\mathrm{CL}}
=
-\frac{1}{N}
\sum_{i=1}^{N}
\log
\frac{
e^{S'_{ii}}
}{
\sum_{j=1}^{N}e^{S'_{ij}}
}.
\end{equation}

Conventional contrastive learning methods generally treat all negative samples equally, which limits their ability to focus on acoustically confusing speaker pairs. To improve discriminative representation learning, we introduce hard negative mining (HNM), which selects the most challenging negative sample for each anchor:
\begin{equation}
S_i^{\mathrm{hard}}
=
\max_{j\neq i}
S_{ij}.
\end{equation}

The corresponding HNM loss is formulated as
\begin{equation}
\mathcal{L}_{\mathrm{HNM}}
=
\frac{1}{N}
\sum_{i=1}^{N}
\max
\left(
0,
S_i^{\mathrm{hard}}
-
S_{ii}
+
m
\right),
\end{equation}
where $m$ denotes the margin parameter controlling the separation between positive and negative pairs.

The final optimization objective jointly combines reliability-aware pseudo-label supervision and discriminative contrastive representation learning:
\begin{equation}
\mathcal{L}_{\mathrm{Total}}
=
\mathcal{L}_{\mathrm{DLNC}}
+
\gamma
\mathcal{L}_{\mathrm{HNM}},
\end{equation}
where $\gamma$ controls the contribution of the HNM objective.
\subsection{Offline Pseudo-Label Correction}
The proposed DLNC performs online confidence-aware supervision during model optimization without explicitly modifying pseudo-label assignments. To further refine pseudo-label quality across successive training iterations, we additionally adopt the offline pseudo-label correction strategy proposed in [36]. Specifically, each unlabeled utterance is divided into
$M$ speech segments and the trained model predicts
a speaker label for each segment individually.
Let $Y_t$ denote the most frequently predicted label. If
\begin{equation}
\frac{I}{M} > 0.5,
\end{equation} 
where $I$ is the occurrence count of $Y_t$,
the pseudo-label is updated to $Y_t$.
Otherwise, the sample is regarded as unreliable and
excluded from the subsequent training iteration. The corrected pseudo-label set is then used to initialize
the next round of iterative optimization.

\section{Experimental Setup}

\subsection{Datasets}

\textbf{Training Set}: 
To evaluate the proposed HCPL framework, experiments were conducted on the VoxCeleb2 dataset [43], which is a large-scale speaker verification corpus collected from YouTube interview videos. VoxCeleb2 contains 1,092,009 utterances from 5,994 speakers and serves as an extension of the original VoxCeleb dataset. In this work, only the audio modality was utilized, and no speaker identity annotations were used during the entire training process.

\textbf{Evaluation Set}: 
The proposed method is evaluated on three standard VoxCeleb1 evaluation protocols [60]: Vox1-O, Vox1-E, and Vox1-H. Vox1-O corresponds to the original VoxCeleb1 test set and contains 37,720 verification trials from 40 speakers. Vox1-E extends the evaluation to the entire VoxCeleb1 dataset, consisting of 581,480 verification trials from 1,251 speakers. Vox1-H is a more challenging evaluation protocol containing 552,536 trials from 1,190 speakers, where all speaker pairs share the same nationality and gender, thereby substantially increasing speaker discrimination difficulty.

Performance is evaluated using Equal Error Rate (EER), where the false acceptance rate equals the false rejection rate.


\textbf{Data Augmentation}: 
To improve robustness under complex acoustic conditions, online data augmentation was employed during training. Specifically, additive noise and reverberation were introduced using the MUSAN [61] and RIR [62] datasets.

\subsection{PTM Selection}

Speaker-related information is unevenly distributed across the
Transformer layers of speech pre-trained models. We therefore conduct
a layer-wise evaluation on the Vox1-O test set to identify the
representations that provide the strongest speaker discrimination.
A lower EER indicates that the corresponding layer preserves more
speaker-discriminative information.

As illustrated in Fig.~\ref{fig:Transformer_layers}, WavLM achieves
lower EERs than the other evaluated PTMs across most Transformer
layers. In particular, the 11th layer of WavLM yields the best
performance, while the 9th layer of UniSpeech-SAT achieves the
second-best result. These two representations are therefore selected
as the feature extractors for initial pseudo-label generation.

In addition to their strong individual performance, WavLM and
UniSpeech-SAT employ different pre-training objectives and
speaker-related modeling mechanisms. Their representations may
therefore induce non-identical clustering partitions and assignment
errors. This potential diversity motivates the use of cross-model
assignment consistency as a reliability cue in the proposed HCPL
framework.

\begin{figure}[!t]
\centering
\includegraphics[width=3.5in]{PTM_Select.png}
\caption{\centering{EER\% performance on Vox1-O across different Transformer layers.}}
\label{fig:Transformer_layers}
\end{figure}

\subsection{Clustering Analysis}

To select an appropriate clustering algorithm for pseudo-label
generation, we evaluate K-means, Agglomerative Hierarchical
Clustering (AHC) [63], Leiden [64], and Infomap on
VoxCeleb2. Speaker embeddings are extracted from the 11th layer of
WavLM and the 9th layer of UniSpeech-SAT, which provide the strongest
speaker discrimination in the preceding layer-wise analysis.

For K-means and AHC, the number of clusters is set to 6,000, which
approximately matches the number of speakers in VoxCeleb2. This
setting is used only for comparative evaluation because these two
algorithms require a predefined number of clusters. In contrast,
Infomap and Leiden do not use the ground-truth speaker count during
clustering. Ground-truth speaker identities are used only to compute
the clustering metrics and are not involved in pseudo-label
generation or model optimization.Clustering performance is evaluated using Precision, Recall, F1-score, and Normalized Mutual Information (NMI). 

As summarized in Table~\ref{tab:clustering}, Leiden achieves the
highest precision for both PTM feature sets, but its substantially
lower recall indicates that it tends to produce highly fragmented
clusters. By contrast, Infomap achieves the highest F1-score and NMI
for both WavLM and UniSpeech-SAT embeddings. This result indicates a
more favorable balance between cluster purity and cluster
completeness. Infomap is therefore adopted as the clustering
algorithm in the subsequent experiments.


\begin{table}[!t] 
    \centering
    \caption{Clustering quality of pseudo-label sets generated during the first iteration.} % 表格标题
    \setlength{\tabcolsep}{4pt} % 微调列间距（可选）
    \begin{tabular}{c c c c c c} % 补全竖线，更规范
        \hline
        PTM & Method & Precision$\uparrow$ & Recall$\uparrow$ & F1-score$\uparrow$ & NMI$\uparrow$ \\
        \hline
        \multirow{4}{*}{WavLM}
        & K-means  & 0.303 & 0.209 & 0.247  & 0.665  \\ 
        & AHC  & 0.442 & 0.346 & 0.388  & 0.776  \\
        & Leiden  & \textbf{0.887} & 0.237 & 0.374  & 0.859  \\
        & Infomap & 0.795 & \textbf{0.426} & \textbf{0.555} &\textbf{0.871}  \\
        \hline
        \multirow{4}{*}{UniSpeech-SAT}
        & K-means  & 0.287 & 0.200 & 0.236  & 0.651  \\ 
        & AHC  & 0.442 & 0.347 & 0.389  & 0.787  \\
        & Leiden  & \textbf{0.901} & 0.173 & 0.291  & 0.861  \\
        & Infomap  & 0.697 & \textbf{0.364} & \textbf{0.579}&\textbf{0.872}  \\
        \hline
    \end{tabular}
    \label{tab:clustering}
\end{table}



\subsection{First-Round Pseudo-Label Analysis}
\label{sec:first_round_analysis}

To provide a more comprehensive evaluation of pseudo-label quality,
we additionally report Clustering Accuracy (Cluster-Acc), which
measures the maximum sample-level agreement between pseudo-label
assignments and ground-truth speaker identities after optimal label
alignment. The ground-truth identities are used only for offline
evaluation and are not involved in pseudo-label generation or model
training.

Cluster-Acc is defined as
\begin{equation}
\mathrm{Cluster\mbox{-}Acc}
=
\frac{
\max_{\pi}
\sum_{(i,j)\in\pi} C(i,j)
}{
M
},
\end{equation}
where $C(i,j)$ denotes the number of samples shared by the $i$-th
pseudo-label cluster and the $j$-th ground-truth speaker class,
$\pi$ denotes a feasible one-to-one mapping between pseudo-label
clusters and ground-truth classes, and $M$ denotes the total number
of evaluated samples. The optimal mapping is obtained using the
Hungarian algorithm [55].

Using Infomap clustering, we first generate two initial pseudo-label
sets, denoted as $Y_A$ and $Y_B$, from WavLM and UniSpeech-SAT
embeddings, respectively. The two clustering outputs are aligned, and
samples receiving consistent assignments across the two PTMs are
retained by HCF. The resulting pseudo-label set is denoted as
$Y_{\mathrm{hc}}$.

As summarized in
Table~\ref{tab:First-round pseudo-label set metrics}, HCF retains
755,517 out of 1,092,009 training utterances, corresponding to a
sample retention rate of 69.19\%. Compared with the two original
pseudo-label sets, $Y_{\mathrm{hc}}$ achieves consistently higher
clustering quality. Its F1-score increases from 0.555 and 0.479 to
0.692, its NMI increases from 0.871 and 0.872 to 0.942, and its
Cluster-Acc increases from 0.569 and 0.519 to 0.677, relative to
$Y_A$ and $Y_B$, respectively.

Although HCF removes 30.81\% of the utterances from
pseudo-label-supervised training, the reduction in sample quantity is
accompanied by a substantial improvement in supervision quality.
Samples with inconsistent assignments across the two clustering
outputs are more likely to lie near ambiguous clustering boundaries
or to be affected by model-specific partitioning errors. In
pseudo-label-based learning, an incorrectly assigned sample may be
more harmful than the temporary absence of a pseudo-label because it
produces gradients toward an incorrect speaker class. HCF therefore
prioritizes a smaller but more coherent supervised subset during the
initial training stage.

The improvement in F1-score, NMI, and Cluster-Acc indicates that,
for the selected WavLM and UniSpeech-SAT representations,
cross-model consistency provides an effective reliability cue for
identifying a cleaner pseudo-label subset. This result validates the
agreement-based filtering criterion for the selected PTMs, without
requiring the assumption that all discarded samples are necessarily
incorrect.

Pseudo-label merging is subsequently applied to
$Y_{\mathrm{hc}}$, producing the final pseudo-label set
$\tilde{Y}$. Compared with $Y_{\mathrm{hc}}$, $\tilde{Y}$ improves
the F1-score from 0.692 to 0.702, the NMI from 0.942 to 0.944, and
the Cluster-Acc from 0.677 to 0.683. Meanwhile, the number of
pseudo-label classes decreases from 18,180 to 17,819, while the
number of retained samples remains unchanged.

The simultaneous reduction in class number and improvement in the
three clustering metrics suggests that PLM consolidates a subset of
highly similar clusters without degrading the measured clustering
quality. This behavior is consistent with the intended role of PLM
in alleviating category fragmentation. Nevertheless, the remaining
number of pseudo-label classes is still larger than the number of
underlying speakers, indicating that PLM reduces rather than
completely eliminates over-segmentation during the initial
iteration.

Based on the above analysis, $\tilde{Y}$ is adopted as the initial
pseudo-label set for downstream unsupervised speaker verification
training. The utterances excluded by HCF are not permanently
discarded. Instead, they are further utilized through the
self-supervised contrastive objective with hard negative mining.


\begin{table}[!t] % table* 表示跨双栏
    \centering
    \caption{First-round pseudo-label set metrics} % 表格标题
    \setlength{\tabcolsep}{4pt} % 微调列间距（可选）
    \begin{tabular}{c c c c c c c} % 补全竖线，更规范
        \hline
        Label set  & F1-score$\uparrow$ & NMI$\uparrow$ & Cluster-Acc$\uparrow$ & Class-Num & Sample-Num \\
        \hline
        $Y_A$    & 0.555 & 0.871 & 0.569  & 20385  &1092009  \\    
        $Y_B$      & 0.479  & 0.872 & 0.519 & 21296 & 1092009  \\         
        $Y_{\mathrm{hc}}$      & 0.692 & 0.942 & 0.677  & 18180   & 755517   \\
        $\tilde{Y}$     & \textbf{0.702} & \textbf{0.944} & \textbf{0.683} &  17819  &  755517     \\
        \hline
    \end{tabular}
    \label{tab:First-round pseudo-label set metrics}
\end{table}
\subsection{Implementation Details}

For downstream unsupervised speaker verification training, WavLM is
selected as the PTM backbone because it provides the strongest
speaker discrimination in the layer-wise evaluation. The downstream
speaker encoder is based on ECAPA-TDNN with a channel size of 1024.

Training utterances are segmented into non-overlapping 2-second
speech chunks. The batch size for pseudo-label-supervised training is
set to 512. AdamW is used for optimization together with a
triangular cyclic learning-rate schedule. The training procedure
contains two stages:

\begin{enumerate}
    \item The WavLM parameters are first frozen, and only the
    downstream ECAPA-TDNN is optimized for 10 epochs. The learning
    rate varies from $1\times10^{-8}$ to $5\times10^{-4}$ under a
    triangular cyclic schedule with a cycle step size of 5 epochs.

    \item The WavLM parameters are subsequently unfrozen, and the
    entire model is jointly fine-tuned for another 10 epochs. The
    maximum learning rate is reduced to $5\times10^{-5}$ during this
    stage to stabilize joint optimization.
\end{enumerate}

% The PLM scaling coefficient is fixed to $\mu=0.97$ in all
% experiments and remains unchanged throughout training. The EMA
% momentum in DLNC is set to $v=[VALUE]$, and the remaining
% AAM-Softmax parameters are set to $s=[VALUE]$ and
% $\alpha=[VALUE]$.

For contrastive learning with hard negative mining, the mini-batch
size is set to 16. The relatively small batch size is intended to
reduce, although not completely eliminate, the probability that
multiple utterances from the same speaker appear in the same batch
and are incorrectly treated as negative pairs. The margin parameter
in $\mathcal{L}_{\mathrm{HNM}}$ is fixed to 0.2. During inference,
cosine similarity scoring is used without additional score
normalization.
\begin{table*}[t] % table* 表示跨双栏
    \centering
   \caption{Comparison of the proposed HCPL method with representative single-stage unsupervised speaker verification frameworks in terms of EER (\%) on the Vox1-O, Vox1-E, and Vox1-H test sets.} % 表格标题
    \setlength{\tabcolsep}{10pt} % 微调列间距（可选）
    \begin{tabular}{l l l  c c c} % 补全竖线，更规范
        \hline
        Category & Method  & Model   & Vox1-O & Vox1-E & Vox1-H  \\
        \hline
        \multirow{10}{*}{Contrastive}
        % &AP + l-mix  & ECAPA-TDNN & 10.49 &  &    \\
        % &AP + AAT  & ResNet34 & 8.65 &  &    \\
        &SimCLR + MSE[10] & ResNet34 & 8.28 &  &    \\
        &MoCo + WavAug[11]& TDNN & 8.23 &-  &-    \\
        &SimCLR + Margin[12]& ResNet34 & 7.85 & - &  -  \\
        &SimCLR + AAT[18]& ECAPA-TDNN & 7.36 & - & -   \\
        &C3-MoCo[13]& ECAPA-TDNN & 6.40 &-  & -   \\
        &MCL-DPP[19]& ECAPA-TDNN & 2.89 & 3.17 & 6.27   \\
        &SimCLR-SSPS[14]& ECAPA-TDNN & 2.57 &3.11  & 5.56   \\
        \hline
        \multirow{9}{*}{Self-distillation}
        &DINO + Cosine loss[22]&ResNet34&  6.16 & - & -   \\
        &DINO + Curriculum[23]&ECAPA-TDNN&  4.47 & - & -   \\
        &CA-DINO[21]&ECAPA-TDNN& 3.59 & 3.85 &  6.92  \\
        &RDINO[24]&ECAPA-TDNN& 3.29 & - & -   \\
        &EBCA-DINO[25]&ECAPA-TDNN& 2.94 &2.99  &  5.14  \\
        &DINO-SSPS[14]&ECAPA-TDNN& 2.53 &2.55  &  4.93  \\
        &DINO + Aug[26]&ECAPA-TDNN& 2.51 &2.47  &  4.79  \\
        &C3-DINO[13]&ECAPA-TDNN& 2.50& - &-    \\
        \hline
        \multirow{2}{*}{PTM}
        &Sub-PTM[36]&WavLM\&ECAPA-TDNN& 2.55& - &-    \\
        &HCPL(ours) & WavLM\&ECAPA-TDNN  & \textbf{2.33} & \textbf{2.37} & \textbf{4.28}   \\
        \hline
    \end{tabular}
    \label{tab:single-stage comparison}
\end{table*}
\FloatBarrier 
\subsection{Verification Performance}

Table~\ref{tab:single-stage comparison} compares the proposed HCPL framework with representative single-stage unsupervised speaker verification methods. After only one training iteration, HCPL achieves EERs of 2.33\%, 2.37\%, and 4.28\% on Vox1-O, Vox1-E, and Vox1-H, respectively, outperforming the compared single-stage approaches across all three evaluation protocols. These results indicate that the proposed framework provides a competitive speaker verification model before iterative pseudo-label refinement is performed.

A notable observation is that the performance improvement is consistently maintained across the three evaluation protocols despite their different levels of difficulty. While Vox1-O represents the standard evaluation protocol, Vox1-E introduces substantially more verification trials, and Vox1-H further increases task difficulty by restricting trial pairs to speakers with the same nationality and gender. The consistent reduction in EER across these increasingly challenging settings suggests that the proposed framework improves speaker discrimination under both conventional and difficult verification conditions, rather than benefiting only a specific evaluation protocol. Nevertheless, this observation is limited to the VoxCeleb benchmarks evaluated in this work and should not be interpreted as evidence of generalization to other datasets or application scenarios.

The strong first-round performance is consistent with the pseudo-label quality analysis presented in Table~\ref{tab:First-round pseudo-label set metrics}. One possible explanation is that the proposed framework improves supervision reliability before speaker encoder optimization begins. Specifically, HCF removes pseudo-label assignments that are inconsistent across different speech pre-trained models, thereby reducing the influence of unreliable clustering results. PLM further alleviates category fragmentation among highly similar clusters, providing a more coherent supervisory signal for downstream optimization. During model training, the proposed DLNC objective may further reduce the influence of residual pseudo-label noise through confidence-aware supervision, whereas HNM enables the utterances excluded from pseudo-label-guided optimization to continue contributing to representation learning. Since these components operate at different stages of the framework, their combined effect provides a plausible explanation for the superior first-round performance.

Table~\ref{tab:multi-stage comparison} further compares the final performance of representative iterative unsupervised speaker verification frameworks. After three refinement iterations, the proposed method achieves EERs of 1.01\%, 1.36\%, and 2.53\% on Vox1-O, Vox1-E, and Vox1-H, respectively, obtaining the best performance among the compared unsupervised systems. More importantly, the proposed framework reaches its best reported performance after only three iterations, whereas several competing iterative approaches require considerably more refinement rounds to obtain their final results.

The improvement from the first-round results to the final performance suggests that reliable pseudo-label initialization and subsequent iterative refinement complement each other throughout optimization. Rather than relying solely on repeated pseudo-label updates, the proposed framework starts from a cleaner supervisory subset and progressively refines both the speaker representations and pseudo-label assignments during later iterations. This observation is consistent with the primary objective of this work, namely reducing the propagation of unreliable pseudo-labels during iterative unsupervised speaker verification. Although the current experiments do not directly isolate the contribution of each refinement stage to convergence behaviour, the results suggest that improving pseudo-label reliability during the initial stage is associated with more stable iterative optimization and ultimately lower verification error.

\begin{table*}[t] % table* 表示跨双栏
    \centering
    \caption{Comparison of the proposed HCPL method with representative multi-stage unsupervised speaker verification frameworks in terms of EER (\%) on the Vox1-O, Vox1-E, and Vox1-H test sets.} % 表格标题
    \setlength{\tabcolsep}{10pt} % 微调列间距（可选）
    \begin{tabular}{l l l l c c c c} % 补全竖线，更规范
        \hline
        Method & Stage I & Model & Clustering & Iter & Vox1-O & Vox1-E & Vox1-H  \\
        \hline
        Fully-supervised & - & WavLM\&ECAPA-TDNN& -&1   & 0.94 & 1.19 & 2.25   \\  
        \hline
        Duke-1[15] & SimCLR &ResNet34 &K-Means& 5   & 3.45 & 4.02 & 5.57\\ 
        IDlab[16] & MoCo&ECAPA-TDNN &AHC & 7   & 2.10 & - & -   \\ 
        JHU[27] & DINO&Res2Net50 &AHC & 4   & 1.89 & - & -   \\ 
        Duke-2[17] &SimCLR& ResNet34& K-Means &4   & 1.81 & - & -   \\ 
        LGL[18]  & SimCLR&ECAPA-TDNN &K-Means& 5   & 1.66 & 2.18 & 3.76   \\ 
        CA-DINO[21] & DINO &ECAPA-TDNN&K-Means& 3   & 1.58 & 1.88 & 3.29   \\ 
        DLG-LC[28]  & DINO &ECAPA-TDNN &K-Means& 5   & 1.47 & 1.78 & 3.19   \\ 
        MCL-DPP[19]& SimCLR &ECAPA-TDNN &K-Means& 4   & 1.44 & 1.77 & 3.27   \\ 
        AT + HT[29]  & DINO &ECAPA-TDNN &K-Means& 5   & 1.35 & - & -   \\ 
        BDS-BPLC[30] & DINO &ECAPA-TDNN&K-Means& 5   & 1.33 & 1.56 & 2.78   \\ 

        Co-Meta[20] & SimCLR &ECAPA-TDNN&K-Means& 7   & 1.27 & 1.82 & 3.31   \\ 
        Sub-PTM[36] & WavLM & WavLM\&ECAPA-TDNN&Infomap & 5   & 1.25 & - & -   \\ 
        IPL[65] &i-vector &MFA-Conformer&  AHC &8   & 1.06 & 2.16 & 3.61   \\ 
        HCPL(ours) &WavLM\&UniSpeech-SAT & WavLM\&ECAPA-TDNN&Infomap& 3   & \textbf{1.01} & \textbf{1.36} & \textbf{2.53}   \\
        \hline
    \end{tabular}
    \label{tab:multi-stage comparison}
\end{table*}
\FloatBarrier

\begin{table}[!t] % table* 表示跨双栏
    \centering
    \caption{ Detailed results for each iteration. All metrics are reported in B/A format, where B denotes the value before online label correction and A after offline label correction.} % 表格标题
    \setlength{\tabcolsep}{3.6pt} % 微调列间距（可选）
    \begin{tabular}{c c c c} % 补全竖线，更规范
        \hline
        Method & Iter1 & Iter2 & Iter3   \\
        \hline
        F1-score$\uparrow$  & 0.702 / 0.719 & 0.826 / 0.865 & 0.883 / 0.907 \\  
        NMI$\uparrow$  & 0.944 / 0.945 & 0.962 / 0.973 & 0.976 / 0.981 \\
        Cluster-Acc$\uparrow$  & 0.683 / 0.707 & 0.868 / 0.879 & 0.916 / 0.924 \\
        Class-Num  & 17819 / 14628 & 11501 / 9447 & 11119 / 8820 \\
        Sample-Num  & 755517 / 991261 & 1092009 / 1049500 & 1092009 / 1068949 \\
        \hline
    \end{tabular}
    \label{tab:iteration}
\end{table}
\FloatBarrier

\subsection{Iterative Pseudo-Label Refinement}

Table~\ref{tab:iteration} summarizes the evolution of pseudo-label quality and training data coverage during iterative optimization. Across the three refinement iterations, F1-score, NMI, and Cluster-Acc all increase consistently, indicating that the pseudo-label assignments become progressively more coherent and better aligned with the underlying speaker structure. In addition, the improvements are observed both before and after offline pseudo-label correction, suggesting that the proposed framework continuously enhances pseudo-label quality throughout the iterative learning process rather than relying on a single refinement stage.

A second notable trend is the progressive recovery of training samples. During the first iteration, HCF deliberately retains only the subset of samples with consistent cross-model assignments to suppress the influence of unreliable supervision. As iterative optimization proceeds, the number of samples participating in pseudo-label-guided training increases substantially after offline correction. This observation suggests that the cleaner supervisory signal obtained during the early stage enables the speaker encoder to produce increasingly reliable predictions, allowing a larger proportion of previously excluded utterances to be incorporated into subsequent training iterations. Consequently, HCF functions as a conservative initialization strategy for controlling early-stage error propagation rather than permanently discarding difficult samples.

The number of pseudo-label classes decreases considerably throughout the refinement process. This trend is consistent with a gradual reduction in category fragmentation, where multiple pseudo-label clusters corresponding to the same speaker are progressively consolidated during iterative optimization. Although the final number of pseudo-label classes remains higher than the actual number of speakers in VoxCeleb2, this residual over-segmentation does not prevent the proposed framework from achieving strong verification performance. One possible explanation is that speaker verification primarily depends on the discriminative geometry of the embedding space rather than an exact one-to-one correspondence between pseudo-label classes and speaker identities. As long as utterances from the same speaker remain sufficiently close while different speakers remain well separated in the embedding space, moderate over-segmentation may still support effective verification performance.

Another noteworthy observation is that the improvements gradually diminish in the later iterations. Compared with the substantial gains obtained during the first two rounds, the changes in clustering quality become relatively small after the third refinement iteration. This trend suggests that the speaker representations and pseudo-label assignments progressively approach a stable state under the proposed iterative optimization framework. Since the model architecture and optimization strategy remain unchanged during refinement, the observed improvement is more likely associated with progressively improved supervision quality than with increased model capacity.

Overall, the iterative refinement results are consistent with the primary objective of the proposed framework. Instead of repeatedly propagating unreliable pseudo-labels, the framework first establishes a relatively reliable supervisory subset and then progressively expands the effective training data while improving pseudo-label consistency. Although the present experiments do not directly quantify the contribution of each refinement component to convergence behaviour, the observed trends suggest that reliable pseudo-label initialization is associated with more stable iterative optimization and improved pseudo-label quality on the evaluated benchmarks.










\begin{table}[!t] % table* 表示跨双栏
    \centering
    \caption{Component-wise ablation results on Vox1-O in terms of EER (\%).} % 表格标题
    \setlength{\tabcolsep}{6.5pt} % 微调列间距（可选）
    \begin{tabular}{c c c c c c c} % 补全竖线，更规范
        \hline
        $\mathcal{L}_{PL}$ & HCF & PLM & $\mathcal{L}_{DLNC}$  & $\mathcal{L}_{CL}$ & $\mathcal{L}_{HNM}$ & EER (\%)\\
        \hline
        $\checkmark$ & & & & & & 2.776  \\  
        $\checkmark$ & $\checkmark$ & & & & & 2.431\\
        $\checkmark$ & $\checkmark$  & $\checkmark$ & & & & 2.415\\
        &$\checkmark$ & $\checkmark$  & $\checkmark$ & & & 2.384\\
        &$\checkmark$ & $\checkmark$  & $\checkmark$ &$\checkmark$ & & 2.369\\
        &$\checkmark$ & $\checkmark$  & $\checkmark$ & &$\checkmark$ & 2.334\\
        \hline
    \end{tabular}
    \label{tab:ablation}
\end{table}
\FloatBarrier





\subsection{Ablation Study and Contrastive Loss Analysis}

Table~\ref{tab:ablation} presents the ablation results on Vox1-O after the first training iteration. Evaluating all variants under the same training condition allows the contribution of each component to be assessed before iterative refinement further improves pseudo-label quality. As additional modules are incorporated, the EER decreases consistently, suggesting that the proposed framework benefits from the complementary effects of pseudo-label refinement and noise-robust optimization rather than relying on a single component.

Among the pseudo-label generation modules, HCF provides the largest performance improvement, indicating that improving the reliability of the initial supervisory signal is particularly important during early-stage optimization. PLM further reduces the EER, suggesting that alleviating category fragmentation provides additional benefits even after unreliable pseudo-label assignments have been filtered.
Replacing the conventional pseudo-label loss with DLNC leads to a further improvement in performance. One possible explanation is that confidence-aware soft correction reduces the influence of residual pseudo-label noise while avoiding the instability associated with directly replacing pseudo-labels using model predictions.
Both contrastive learning strategies further improve verification performance, indicating that the samples excluded from pseudo-label-guided training still provide useful supervisory information. Compared with the conventional contrastive objective, HNM achieves a lower EER, suggesting that emphasizing acoustically confusing negative pairs produces more discriminative speaker representations.

Overall, the ablation results suggest that the proposed modules address different stages of the training pipeline. HCF and PLM improve pseudo-label quality before optimization, whereas DLNC and HNM enhance robustness during representation learning. Their cumulative improvement indicates that reliable pseudo-label generation and noise-aware optimization jointly contribute to the final verification performance.


\begin{figure}[!t]
\centering
\includegraphics[width=3in]{Contra_loss.png}
\caption{\centering{Effect of contrastive loss weight $\gamma$ on Vox1-O EER after the first training iteration.}}
\label{fig:contrastive_loss}
\end{figure}

Fig.~\ref{fig:contrastive_loss} further illustrates the influence of the contrastive loss weight $\gamma$ on the performance of both $L_{\mathrm{CL}}$ and $L_{\mathrm{HNM}}$ after the first training iteration. For both objectives, the best performance is achieved at $\gamma=0.2$, whereas both smaller and larger values lead to higher EERs. This trend suggests that an appropriate balance between pseudo-label supervision and contrastive learning is important during the early stage of optimization.

When $\gamma=0.1$, the contribution of the contrastive objective appears to be limited, resulting in relatively modest performance gains. In contrast, increasing $\gamma$ to 0.2 provides the largest improvement, suggesting that moderate contrastive supervision effectively complements pseudo-label-guided optimization. However, further increasing $\gamma$ to 0.5 or 1.0 consistently degrades performance for both contrastive objectives. One possible explanation is that excessive emphasis on contrastive learning reduces the relative influence of pseudo-label supervision, making it more difficult for the model to fully exploit the reliable supervisory information provided by the proposed HCPL framework.

 Overall, these results suggest that both the design of the contrastive objective and its relative optimization weight influence the effectiveness of the proposed reliability-aware learning framework.



\section{Conclusion}

This paper investigated pseudo-label reliability in iterative unsupervised speaker verification and proposed a High-Confidence Pseudo-Label (HCPL) framework to improve supervision quality during both pseudo-label generation and model optimization. By combining cross-model consistency filtering, pseudo-label merging, dynamic label noise correction, and hard negative mining, the proposed framework addresses unreliable pseudo-label assignments, category fragmentation, and residual label noise within a unified iterative learning framework.

Experimental results on the VoxCeleb benchmarks show that the proposed framework consistently improves speaker verification performance while reaching its best performance after only three refinement iterations. The results further suggest that improving pseudo-label reliability during the early stage of iterative learning is associated with more stable pseudo-label refinement and more effective representation learning throughout subsequent optimization. These findings highlight the importance of supervision quality, in addition to representation learning, for iterative unsupervised speaker verification.

Future work will investigate the applicability of the proposed framework to cross-lingual and cross-domain speaker verification scenarios, where pseudo-label reliability may be further challenged by distribution shifts. In addition, more adaptive pseudo-label refinement strategies and efficient self-supervised optimization methods will be explored to improve robustness in large-scale and low-resource acoustic environments.

% \begin{thebibliography}{1}
% \bibliographystyle{IEEEtran}
\nocite{*}
\bibliographystyle{unsrt}

\bibliography{ref}

% \end{thebibliography}


% \newpage

% \section{Biography Section}
% If you have an EPS/PDF photo (graphicx package needed), extra braces are
%  needed around the contents of the optional argument to biography to prevent
%  the LaTeX parser from getting confused when it sees the complicated
%  $\backslash${\tt{includegraphics}} command within an optional argument. (You can create
%  your own custom macro containing the $\backslash${\tt{includegraphics}} command to make things
%  simpler here.)
 
% \vspace{11pt}

% \bf{If you include a photo:}\vspace{-33pt}
% \begin{IEEEbiography}[{\includegraphics[width=1in,height=1.25in,clip,keepaspectratio]{fig1}}]{Michael Shell}
% Use $\backslash${\tt{begin\{IEEEbiography\}}} and then for the 1st argument use $\backslash${\tt{includegraphics}} to declare and link the author photo.
% Use the author name as the 3rd argument followed by the biography text.
% \end{IEEEbiography}

% \vspace{11pt}

% \bf{If you will not include a photo:}\vspace{-33pt}
% \begin{IEEEbiographynophoto}{John Doe}
% Use $\backslash${\tt{begin\{IEEEbiographynophoto\}}} and the author name as the argument followed by the biography text.
% \end{IEEEbiographynophoto}




\vfill

\end{document}


